\beginsong{Mantas de cartón}[]

    \capo{3}

    \beginverse
        \[Sol]Sale el \[Re]sol, y la \[Mim]luna se esconde de\[Do]prisa.
        \[Sol]Hay niños que \[Re]hoy ama\[Mim]necen con manos \[Do]tendidas.
        \[Sol]En las a\[Re]ceras, la ino\[Mim]cencia vestida de har\[Do]apos,
        \[Sol]mendiga a \[Re]ciegas, lo que al \[Mim]mundo le sobra en las \[Do]manos
    \endverse

    \beginverse
        \[Sol]Sale el \[Re]sol, que ilu\[Mim]mina a ricos y a \[Do]pobres.
        \[Sol]Pregunto \[Re]yo, si el \[Mim]derecho a vivir levan\[Do]tados,
        \[Sol]no será el \[Re]mismo para \[Mim]todos los que pade\[Do]cemos el sol, 
        \[Sol]que nadie ten\[Re]ga que vi\[Mim]vir arras\[Do]trado
    \endverse

    \beginchorus
        Y ese \[Mim]niño en la \[Re]calle, que nun\[Sol]ca ha tenido a \[Mim]nadie,
        que le \[Do]bese las me\[Re]jillas al dor\[Sol]mir.
        Buena \[Mim]madre del \[Re]hombre, \[Sol]acaricia tu sus \[Mim]noches, 
        que tu \[Do]manto no sean \[Re]mantas de car\[Sol]tón
    \endchorus
\endsong
